\section*{الفصل \ref{ch:b3:cell-cycle-apoptosis} --- التحكم في دورة الخلية والموت الخلوي المبرمج}

\begin{solution}{exo:b3:cell-cycle-apoptosis:1}
G1 وS وG2 وM. فتقدم G1 \omterm{def:b3:cell-cycle-apoptosis:checkpoints}{ونقطة التقييد}: \omterm{def:b3:cell-cycle-apoptosis:cdk}{السيكلين} D--Cdk4/6؛
وG1/S: \omterm{def:b3:cell-cycle-apoptosis:cdk}{السيكلين} E--Cdk2؛ وS: \omterm{def:b3:cell-cycle-apoptosis:cdk}{السيكلين} A--Cdk2؛ وG2/M: \omterm{def:b3:cell-cycle-apoptosis:cdk}{السيكلين} B--Cdk1. و
\omterm{def:b3:cell-cycle-apoptosis:cdk}{معقّد تعزيز الطور الانفصالي} (APC/C) يهدم \omterm{def:b3:cell-cycle-apoptosis:cdk}{السيكلين} B والسكيورين و
ينهي الانقسام الفتيلي.
\end{solution}

\begin{solution}{exo:b3:cell-cycle-apoptosis:2}
هارتويل: طافرات \emph{cdc} في خميرة التبرعم، متوقفةً عند أطوار
محددة، ومفهوم البداية. ونيرس: \emph{cdc2} في خميرة الانشطار
كينازًا يوقّت الانقسام الفتيلي، ونظيره البشري (Cdk1) ينقذ
الخميرة --- أي الحفظ. وهانت: \omterm{def:b3:cell-cycle-apoptosis:cdk}{السيكلين} في بيوض قنفذ البحر، بروتين
يُركَّب خلال الدورة ويُهدم عند كل انقسام.
\end{solution}

\begin{solution}{exo:b3:cell-cycle-apoptosis:3}
\omterm{def:b3:cell-cycle-apoptosis:checkpoints}{نقطة التقييد} (في أواخر G1): عوامل النمو والمغذيات والحجم؛ وتعمل على
\omterm{def:b3:cell-cycle-apoptosis:cdk}{السيكلين} D--Cdk4/6 وعلى Rb. وG2/M: سلامة DNA وتضاعفه؛ فيثبط \omterm{def:b3:dna-repair:ddr}{ATM}/ATR
بروتين Cdc25، فيبقى Cdk1 مفسفرًا خاملًا. وتجميع
المغزل: ارتباط كل قسيم مركزي؛ فالقسيمات غير المرتبطة تصنع
معقّدات Mad2--Cdc20 تثبط APC/C.
\end{solution}

\begin{solution}{exo:b3:cell-cycle-apoptosis:4}
\omterm{def:b3:cell-cycle-apoptosis:apoptosis}{الاستماتة}: تنكمش الخلية، ويتكاثف \omterm{def:b3:chromatin-epigenetics:nucleosome}{الصبغين}، ويُقطع DNA سلّمًا،
وينتفط الغشاء لكنه يبقى محكمًا، ويُكشف \omterm{def:b3:cell-cycle-apoptosis:apoptosis}{الفوسفاتيديل سيرين}، وتُبتلع الجثة،
ولا التهاب. \omterm{def:b3:cell-cycle-apoptosis:apoptosis}{والنخر}: تنتفخ الخلية، ويتمزق الغشاء، و
ترشح المحتويات، ويتحلل DNA عشوائيًا، ويقع التهاب. والمنفّذات:
الكاسبازات، وهي بروتيازات سيستئينية تقطع بعد الأسبارتات.
\end{solution}

\begin{solution}{exo:b3:cell-cycle-apoptosis:5}
الارتفاع من $5$ إلى $30$ بمعدل $1$ في الدقيقة: \qty{25}{min}. والتحلل من
$30$ إلى $5$ بعمر نصف \qty{2}{min}: $\log_{2}(30/5)\times 2 =
\qty{5.2}{min}$. فالدور نحو \qty{30}{min}، والانقسام الفتيلي \qty{17}{\%} منه.
\end{solution}

\begin{solution}{exo:b3:cell-cycle-apoptosis:6}
(أ) لا تستطيع $C$ أن تهبط دون $C_{1}$: فتُحبس في الانقسام الفتيلي على الفرع
المرتفع. (ب) بلا فسفرة مثبطة: تنخفض العتبة $C_{2}$،
فيبدأ الانقسام باكرًا عند حجم صغير (وهو النمط الظاهري ``الصغير'').
(ج) يبقى Cdk1 مفسفرًا: فالفرع المرتفع لا يُبلغ، ويقع توقف في
G2 بخلايا مستطيلة. (د) كما في (ب)، ولا تستطيع \omterm{def:b3:cell-cycle-apoptosis:checkpoints}{نقطة تفتيش} G2/M، التي تعمل
عبر Wee1/Cdc25، أن تمسك الخلية بعد ذلك.
\end{solution}

\begin{solution}{exo:b3:cell-cycle-apoptosis:7}
نواة G1 في هيولى الطور S: مناشئها مرخَّصة و
توفر الهيولى \omterm{def:b3:cell-cycle-apoptosis:cdk}{سيكلين} E/A--Cdk2 فعالًا، فتدخل الطور S فورًا. ونواة
G2 في هيولى الطور S: مناشئها أُطلقت ولم يُعد
ترخيصها (فعوامل الترخيص هدمها CDK أو ثبطها)، فلا
تستطيع التضاعف ثانيةً؛ وتنتظر حتى يبلغ الشريك G2 فيدخلان
الانقسام الفتيلي معًا.
\end{solution}

\begin{solution}{exo:b3:cell-cycle-apoptosis:8}
القسيم المركزي غير المرتبط حفّاز: فهو يحوّل Mad2 إلى انتظامه
الفعال باستمرار، فينتج مثبطًا قابلًا للانتشار لبروتين Cdc20
ينتشر في الخلية ويبقي APC/C مطفأً في كل مكان
--- فقسيم واحد يكفي. ومن دون Mad2 يُنشَّط APC/C
متى شغّله Cdk1، سواء ارتبطت الصبغيات أم لا:
فيبدأ الطور الانفصالي بصبغيات غير مرتبطة، وتكون البنات
مختلة الصيغة، ويموت الكائن (أو الخط الخلوي) من فقد الصبغيات.
\end{solution}

\begin{solution}{exo:b3:cell-cycle-apoptosis:9}
بلا \emph{ced-9}: تموت خلايا كان ينبغي أن تعيش --- أي إماتة جنينية
من كثرة الموت. وبلا \emph{ced-3}: لا يقع أي من الموتات $131$،
فتنجو الخلايا وتتمايز، وتكون الدودة قابلة للحياة. وبلا الاثنتين: لا موت
--- أي النمط الظاهري لطافر \emph{ced-3}. ولأن إزالة المنفّذ تلغي
أثر إزالة الحارس، يعمل CED-9 قبلهما، بإمساك
CED-4/CED-3 خاملين؛ فإذا زال كانا فعالين، إلا إن كانا
غائبين هما أيضًا.
\end{solution}

\begin{solution}{exo:b3:cell-cycle-apoptosis:10}
عند $A = aC$، تكون $\mathrm{d}C/\mathrm{d}t = k_{s} - k_{d}aC^{2}$، وهي
معادلة بمتغير واحد حالتها المستقرة $C^{*} = \sqrt{k_{s}/(k_{d}a)}$
وميلها عندها $-2k_{d}aC^{*} < 0$: فأي انحراف يتلاشى
تلاشيًا رتيبًا، والمتغير الواحد من الرتبة الأولى لا يستطيع التذبذب. و
المنحني الذي على شكل S يعطي فرعين مستقرين يفصل بينهما نطاق
ممنوع، فلا بد للحالة أن تقفز ولا تستطيع أن تستقر؛ أما تأخير صريح في
التغذية الراجعة فكان ليعطي تذبذبات بطريق آخر --- أي وصول المكبح
متأخرًا، وتجاوز الهدف، وهكذا --- كما في كثير من
الإيقاعات الهرمونية.
\end{solution}

\begin{solution}{exo:b3:cell-cycle-apoptosis:11}
كل مستشعر يصل يلتقطه حارس، فلا يحدث شيء
حتى يُشغل الحرّاس $N$ كلهم، عند $t \approx N/r$؛ ثم تكون
المستشعرات التالية حرةً، فتنشّط Bax/Bak، الذي يرتفع تكوّن مسامّه ارتفاعًا حادًا
مع عددها، فتنفتح المتقدرات في دقائق. والمنبّه الأشد
لا يقصّر الزمن إلا عبر $r$؛ وفوق العتبة تكون
النتيجة نفسها. والفينيتوكلاكس يشغل الحرّاس، فيخفض
$N$ الفعالة: فيتقلص الزمن، وتموت الخلية المهيّأة --- القريبة سلفًا من $N$
--- فورًا.
\end{solution}

\begin{solution}{exo:b3:cell-cycle-apoptosis:12}
فائض Bcl-2 يحجب \omterm{def:b3:cell-cycle-apoptosis:pathways}{المسلك الجوهري}: فتنجو الخلايا البائية التي كان ينبغي أن تموت
حين تنتهي إشارات نموها، وتنمو الجمهرة بإخفاق
الموت، بالمعدل البطيء الذي تُصنع به تلك الخلايا --- فتكون وديعة،
حتى تضيف آفة ثانية التكاثر. أما فقد Rb فيزيل
\omterm{def:b3:cell-cycle-apoptosis:checkpoints}{نقطة التقييد}: فتمضي طلائع الشبكية في الدورة
بلا عوامل نمو وتنقسم بأسرع ما يسمح به المحرّك ---
فتكون عدوانية. والورمان هما البرنامجان يخفق كل منهما وحده:
ضابطُ موت وضابطُ انقسام، يكفي فقدُ أي منهما
لورم، والثاني أسرع.
\end{solution}

\begin{solution}{pb:b3:cell-cycle-apoptosis:1}
\textbf{1.} $10^{7}\times 3500/5 = 7\times 10^{9}$ خلية في اليوم، أي نحو
\num{80000} في الثانية.
\textbf{2.} $3500/5 = 700$ لكل خبيئة في اليوم؛ و$22\times 2^{5} = 704$.
\textbf{3.} $80\times 365 \approx \num{29000}$ انقسام؛ ومضروبًا في $\times 0.6
\approx \num{17500}$ طفرة.
\textbf{4.} خلايا العبور تُنبذ في أيام، حاملةً طفراتها
معها؛ أما الخلية الجذعية فتبقى مدى الحياة، فتُحفظ كل
طفرة من طفراتها ويضيف كل انقسام مزيدًا --- فالدوران البطيء
يقلّلها.
\textbf{5.} تفرغ الزغابة في نحو $5$ أيام هي زمن عبورها
السوي؛ وتستغرق إعادة البناء يوم تعافٍ، وخمسة انقسامات كل
\qty{12}{h}، والهجرة صعودًا في الزغابة --- أي أربعة إلى خمسة أيام، يكون
الحاجز فيها عاريًا: فانقسامات العبور السريعة في الأمعاء تجعلها
ضحية باكرة للإشعاع.
\textbf{6.} الجثة المستميتة المحكمة لا تطلق شيئًا من محتوياتها في لمعة
مليئة بالبكتيريا ولا تثير التهابًا؛ أما \omterm{def:b3:cell-cycle-apoptosis:apoptosis}{النخر} فكان ليسكب
الإنزيمات والإشارات ويلهب المخاطية عند كل خلية تُنبذ.
\textbf{7.} $(40 - 8)/2 = \qty{16}{h}$.
\textbf{8.} $\log_{2}(40/8)\times \qty{10}{min} = 2.3\times 10 =
\qty{23}{min}$.
\textbf{9.} الدور نحو \qty{16.4}{h}؛ وCdk1 فعال \qty{2.3}{\%} من
الزمن.
\textbf{10.} ليس $k_{s}$ وحدها بل الزمن الذي يُمسك عند العتبات:
فالخلية الجسدية تنتظر عند \omterm{def:b3:cell-cycle-apoptosis:checkpoints}{نقطة التقييد} عوامل النمو وعند
G2/M سلامةَ DNA، وتنتظر الخلية الجذعية أطول من خلية العبور. و
بيضة الضفدع لا نقاط تفتيش لها وهيولاها مجهزة بكل شيء، فلا
يحدد دورها إلا المذبذب العاري.
\textbf{11.} على الفرع المنخفض عند مستوى \omterm{def:b3:cell-cycle-apoptosis:cdk}{سيكلين} فوق
$C_{2}$ السوي: فقد أزاحت \omterm{def:b3:cell-cycle-apoptosis:checkpoints}{نقطة التفتيش} منحني S إلى اليمين
بتثبيط Cdc25، فلا تقع القفزة. وعند الإصلاح يُطلق Cdc25،
فتهبط العتبة دون \omterm{def:b3:cell-cycle-apoptosis:cdk}{السيكلين} المتراكم، وتدخل
الخلية الانقسام الفتيلي فورًا --- ولهذا تدخل الخلايا المحررة من \omterm{def:b3:cell-cycle-apoptosis:checkpoints}{نقطة التفتيش}
الانقسامَ متزامنةً.
\textbf{12.} لا يُحلَّل \omterm{def:b3:cell-cycle-apoptosis:cdk}{السيكلين} B ولا السكيورين: فتبقى الخلية في
الطور الاستوائي والكوهيزين سليم وCdk1 فعال، إلى ما لا نهاية؛ فتموت
هناك أو تنزلق في النهاية خارجه بجينوم غير مفروز.
\textbf{13.} $10^{11}$ انقسام في اليوم.
\textbf{14.} مع \omterm{def:b3:cell-cycle-apoptosis:checkpoints}{نقطة التفتيش} $10^{11}/5\times 10^{4} = 2\times 10^{6}$
بنت مختلة الصيغة في اليوم؛ ومن دونها $2\times 10^{9}$.
\textbf{15.} $2\times 10^{4}$ في اليوم؛ وعلى مدى $80$ سنة $6\times 10^{8}$.
\textbf{16.} $10^{9}/50 = 2\times 10^{7}$ سوء فرز في اليوم: أي إن الورم
يجرب كل يوم عشرين مليون نمط صبغي جديد، ويبقي الانتقاء
ما ينمو أسرع أو يقاوم دواءً.
\textbf{17.} إذا أساء كل انقسام الفرز، لم تحفظ أي سلالة
خلوية جنينية جينومًا سوي الصيغة ومات الجنين. أما الفقد الجزئي فيعطي
اختلال صيغة عارضًا في خلايا تنجو منه (ولا سيما بلا
p53)، أي لااستقرارًا صبغيًا يوفر التنوع لورم
من غير أن يقتل الكائن.
\textbf{18.} خطأ الانقسام المنصّف يضع الصبغي الزائد في العروس،
فتكون اللاقحة وكل خلية منحدرة منها ثلاثية الصبغي؛ أما الخطأ
الفتيلي فيقع في خلية جسدية واحدة ولا ترثه إلا نسيلتها.
\textbf{19.} $10^{11}\times \qty{1}{ng} = \qty{100}{g}$ في اليوم، و
\qty{36}{kg} في السنة --- أي نحو نصف كتلة الجسم كل سنة.
\textbf{20.} \qty{20}{g} من البروتين في اليوم، أي نحو ثلث المدخول
الغذائي وبضعة في المئة من دوران بروتين الجسم الكلي.
\textbf{21.} $10^{11}$ جثة بمعدل $24$ لكل بلعمي في اليوم: أي $4\times
10^{9}$ بلعمي بدوام كامل، أو \qty{4}{\%} من وقت البلعميات البالغة
$10^{11}$.
\textbf{22.} الجثة المتحللة تطلق بروتيازات وDNA وATP وغيرها من
إشارات الإنذار التي تسبب الالتهاب وقد تثير مناعة ذاتية
تجاه مستضدات نووية؛ أما الجثة السليمة فتعلن عن نفسها بجزيء
\omterm{def:b3:cell-cycle-apoptosis:apoptosis}{الفوسفاتيديل سيرين} على وريقتها الخارجية (وبفقدها إشارات ``لا تأكلني'')،
وتجذب البلعميات بنوكليوتيدات مطلَقة.
\textbf{23.} \omterm{def:b3:cell-cycle-apoptosis:pathways}{المسلك الجوهري} محجوب عند المتقدرة. و
الطريق الخارجي ما يزال يعمل حيث ينشّط الكاسباز-8 الكاسبازَ-3
مباشرةً، وإن فُقد تضخيمه عبر Bid. فتتراكم الخلايا
بألا تموت لا بأن تنقسم أسرع --- أي نمو بطيء
--- حتى تضيف آفة لاحقة (وهي \emph{MYC} غالبًا) التكاثر.
\textbf{24.} \omterm{def:b3:cell-cycle-apoptosis:apoptosis}{الاستماتة} تستغرق ساعات، وتحتاج إلى ATP وتمر
بخطوات يمكن حجبها --- مثبطات \omterm{def:b3:cell-cycle-apoptosis:apoptosis}{الكاسباز}، والعتبة
المتقدرية --- فيمكن إنقاذ عصبونات شبه الظل إذا عولجت في الوقت؛
أما عصبونات اللب فتموت من إخفاق الطاقة في دقائق، بلا برنامج
يمكن مقاطعته.
\textbf{25.} دور المذبذب العاري نحو \qty{16}{h}؛ ونحو $2\times
10^{4}$ خلية مختلة الصيغة منقسمة في اليوم في جسم سليم؛ و\qty{100}{g}
من الخلايا معادة التدوير \omterm{def:b3:cell-cycle-apoptosis:apoptosis}{بالاستماتة} كل يوم.
\end{solution}
